\section{Discusión}
\label{sec:discusion}
% Interpretación, comparación con la teoría, hallazgos inesperados, limitaciones.
% Organizado por los mismos tres fenómenos que Resultados.

\subsection{Pérdida de la asociatividad}
\label{subsec:dis-asociatividad}

El contraste entre las Figuras \ref{fig:exp1-int} y \ref{fig:exp1-float} separa dos aritméticas que
la notación matemática no distingue. Con $b$ de tipo \texttt{int}, la identidad $1 + b^{k} - b^{k} = 1$
se traslada al programa sin pérdida porque el entero de Python es exacto y no tiene techo
(sección \ref{subsec:enteros}); las tres asociaciones son entonces tres formas de escribir el mismo
cálculo. Con $b$ de tipo \texttt{float} la identidad deja de valer en cuanto $b^{k}$ alcanza el
tamaño en que el 1 ya no cabe en la mantisa del acumulador. La coincidencia entre el $k^{\ast}$
predicho por la condición $b^{k} \ge 2^{53}$ y la altura de la meseta medida, exacta en seis de las
ocho bases de la Tabla \ref{tab:exp1-umbrales} y con un término de diferencia en las dos restantes,
indica que el efecto no depende de detalles de la implementación
sino del espaciado de $\mathbb{F}$ descrito en \eqref{eq:binary64}. Es el contraejemplo
\eqref{eq:contraejemplo} repetido a lo largo de un barrido.

Conviene precisar qué falla. La potencia $b^{k}$ está bien calculada dentro de lo que el formato
permite; lo que se pierde es el 1, absorbido en la suma $1 + b^{k}$. El término no se anula porque
$b^{k}$ sea inexacto, sino porque la agrupación decide en qué momento el 1 se encuentra con un
número demasiado grande.

Las dos primeras asociaciones dieron resultados idénticos en todas las bases ensayadas, y la razón
es sintáctica antes que numérica: Python evalúa $1 + b^{k} - b^{k}$ de izquierda a derecha, así que
esa expresión y $(1 + b^{k}) - b^{k}$ describen la misma secuencia de operaciones. La consigna pide
comparar tres asociaciones y el lenguaje entrega dos programas distintos. El par idéntico deja
aislado el efecto del agrupamiento: las dos expresiones que se traducen a la misma secuencia
coinciden en todas las bases, y la única que se aparta es la que agrupa distinto.

La tercera asociación se aparta de las otras dos solo en $b = \pm 2.0$, y ahí difiere en un
término: retiene uno más con $b = 2.0$ y uno menos con $b = -2.0$. El mecanismo es la asimetría entre binadas vecinas señalada en la sección
\ref{subsec:punto-flotante}: en $[2^{52}, 2^{53})$ el espaciado vale 1 y en $[2^{53}, 2^{54})$ vale
2, de modo que $1 - 2^{53}$ es representable y $1 + 2^{53}$ no lo es. Restar primero deja el
resultado en la binada de abajo, que tiene el doble de resolución. Que el efecto aparezca
únicamente con $b = \pm 2$ es consecuencia de que solo esa base pisa exactamente el borde $2^{53}$:
con $b = 3$, $5$ o $10$ la potencia salta por encima del borde y las tres asociaciones pierden el
término en el mismo $k$. No alcanza con que la base sea potencia de dos. Con $b = 4$ las potencias
recorren los exponentes pares y saltan de $2^{52}$ a $2^{54}$ sin tocar $2^{53}$, de modo que las
tres asociaciones vuelven a coincidir. La inversión con $b = -2.0$ es la misma asimetría con los
signos cambiados.

Ese mecanismo también da cuenta del régimen $0 \le b \le 1$ de la Tabla
\ref{tab:exp1-b-hasta-uno}. El
fenómeno principal no aparece porque $|b|^{k}$ queda acotado por 1 y nunca se acerca a $2^{53}$; lo
que queda es el redondeo de $1 + b^{k}$, de medio ulp en la binada $[1, 2)$, es decir,
$2^{-53} = 1.11 \times 10^{-16}$, exactamente el valor máximo que registra la tabla para las dos
primeras asociaciones. La tercera, que resta primero y trabaja en la binada de abajo, no muestra
desvío en ninguna de las cinco bases. Que el error de la suma completa sea nulo pese a esos desvíos
es consistente con el mismo espaciado: el acumulador crece hasta 1000, donde el espaciado de
$\mathbb{F}$ vale $2^{-43} \approx 1.14 \times 10^{-13}$. Un desvío de $10^{-16}$ por término cae
por debajo de medio espaciado apenas el acumulador supera unas pocas unidades. El redondeo de la
acumulación descarta el error del término. El mismo mecanismo de absorción que arruina el
experimento con bases grandes lo salva con bases chicas.

El desborde a $\infty$ merece una lectura aparte. Cuando $b^{k}$ excede el mayor representable, el
estándar no interrumpe el cálculo: devuelve $\infty$, y la resta $\infty - \infty$ produce
\texttt{NaN}, que se propaga por el resto de la acumulación. Por eso las curvas de la Figura
\ref{fig:exp1-float} no se desvían: desaparecen. En términos prácticos, el desborde no se manifiesta
como una falla sino como un valor que contamina todo lo que viene después. Detectarlo exige
inspeccionar el resultado; el programa termina igual.

El comportamiento de los enteros de ancho fijo fue el primero de los dos resultados que
contradijeron la expectativa inicial. La
Tabla \ref{tab:exp1-int64} muestra que $b^{k}$ sale del rango de \texttt{int64} muy temprano, en
$k = 19$ para $b = 10$, y que el valor almacenado tiene el signo cambiado en tres de las cuatro
bases. Aun así $a_{1000}$ vale 1000 en las cuatro. La explicación está en la sección
\ref{subsec:enteros}: el entero de 64 bits reduce módulo $2^{64}$ en complemento a dos, y esa
aritmética es la de
un anillo, donde $(1 + x) - x = 1$ vale para todo $x$ sin excepción. El desborde entero no destruye
información, la reduce módulo $2^{64}$, y la resta posterior deshace la reducción. El redondeo de
punto flotante, en cambio, descarta bits de forma irreversible. Esa es la diferencia entre las dos
formas de perder magnitud, y es lo que hace que el mismo experimento revele el problema con
\texttt{float} y lo esconda con \texttt{int64}.

La conclusión que no se sigue de ese dato es que los enteros de ancho fijo sean seguros. La suma que
propone la consigna es justamente una en la que el valor desbordado se cancela contra sí mismo.
Cualquier expresión que use $b^{k}$ sin cancelarlo, por ejemplo una comparación, una división o una
impresión, entrega el valor de la cuarta columna de la Tabla \ref{tab:exp1-int64} y no el de la
tercera. A eso se suma que en C el desborde de un entero con signo es comportamiento indefinido
(ISO, 2018): el mismo programa ni siquiera tiene garantizado el resultado modular que sí produce
NumPy.

La aritmética exacta de Python resuelve el problema pero lo cobra. La Tabla \ref{tab:exp1-costo}
muestra que el tamaño de $b^{N}$ crece linealmente y el tiempo de la corrida completa, de manera
cuadrática, con un cociente contra la corrida en flotantes que crece hasta dos órdenes de magnitud
al ir de $N = 1000$ a $N = 100\,000$. El crecimiento coincide con el argumento de la sección \ref{subsec:enteros}: cada
multiplicación cuesta un tiempo proporcional a la longitud del operando y las longitudes crecen con
$k$. Para $N$ del orden de mil el sobrecosto es tolerable; para $N$ del orden de $10^{5}$, un solo
número ocupa 43.3 KiB y la corrida se vuelve dos órdenes de magnitud más lenta. La elección entre
exactitud y costo depende del tamaño del problema.

\subsection{Dependencia del orden de suma}
\label{subsec:dis-orden}

El primer punto es de nombre. El experimento se plantea sobre la propiedad conmutativa, pero lo que
los datos ponen a prueba es la asociatividad. La suma de punto flotante sí es conmutativa, porque
$\mathrm{fl}(x + y)$ depende del valor exacto de $x + y$ y de una función de redondeo simétrica en
sus argumentos. Los cuatro algoritmos de la Figura \ref{fig:exp2-error} suman exactamente los mismos
números y difieren solo en cómo encadenan las sumas parciales, es decir, en el agrupamiento. La
distinción se enunció en la sección \ref{subsec:asociatividad} y los resultados la sostienen: si el
efecto viniera de la conmutatividad, permutar los índices no debería cambiar nada, y cambia.

La separación de las cuatro curvas en dos grupos sigue lo que predice el análisis de la sección
\ref{subsec:asociatividad}. Los términos $1/(k(k+1))$ decrecen como $k^{-2}$, así que sumar de mayor
a menor lleva el acumulador cerca de 1 en las primeras iteraciones y obliga a los términos tardíos,
de orden $10^{-12}$ en $N = 10^{6}$, a competir con un espaciado local del orden de $10^{-16}$. Cada
uno pierde su parte baja. Sumar de menor a mayor agrupa primero los términos chicos entre sí, donde
el espaciado local es comparable a su tamaño, y esa información sobrevive hasta el final. La suma de
Kahan llega al mismo lugar por otro camino: recupera el residuo de cada redondeo y lo reinyecta, con la cota
$2u + O(Nu^{2})$ de la sección \ref{subsec:kahan}, que no crece con $N$. La Tabla
\ref{tab:exp2-resumen} lo confirma: los máximos de esos dos algoritmos quedan entre
$1.11$ y $1.15 \times 10^{-16}$, es decir, en el nivel de $u$, y no crecen al multiplicar $N$ por
cien.

De ahí sale la lectura práctica, que no es la esperable. Ordenar de menor a mayor módulo alcanza
para igualar a Kahan en esta serie: máximos indistinguibles, los dos al nivel de $u$, y la
diferencia queda en
la cantidad de $N$ con coincidencia bit a bit, 599 contra 862 de mil. Cuando el problema es el
orden, ordenar ya lo resuelve, y Kahan cuesta cuatro operaciones de punto flotante por término en
lugar de una. La ventaja de Kahan aparece cuando el orden no está disponible de antemano: acá los
términos vienen ordenados por construcción y basta recorrer el índice al revés. En una serie cuyas
magnitudes no se conocen sin calcularlas, ordenar exige una pasada previa y memoria para
guardar los términos. Kahan no necesita ninguna de las dos cosas.

Las dos curvas que crecen quedan contenidas por $\sqrt{N}\,u$ de \eqref{eq:cota-raiz} en todo el
rango, lo que respalda el modelo de errores independientes de media nula frente al análisis de peor
caso. La cota conservadora $Nu$ vale $1.11 \times 10^{-10}$ en $N = 10^{6}$ y el error medido es
$5.03 \times 10^{-14}$: sobreestima por más de tres órdenes de magnitud. Suponer que los $N-1$
redondeos se alinean en la misma dirección describe un caso posible y no el caso típico.

La Figura \ref{fig:exp2-error} muestra además una estructura que la teoría anticipa: los errores no
nulos se acomodan en bandas horizontales discretas y no en una nube continua. Como $b_N$ está cerca
de 1, la diferencia entre el resultado y la referencia solo puede tomar valores múltiplos del
espaciado local de $\mathbb{F}$. Lo que se ve en el eje vertical son
cantidades enteras de ulp.

La dispersión de la suma randomizada responde la pregunta de la consigna y agrega un matiz. La
respuesta es que no da siempre lo mismo: 48 resultados distintos en 50 repeticiones, con un error
que va de $2.22 \times 10^{-16}$ a $5.64 \times 10^{-14}$, un factor 254 entre el mejor y el peor
caso, según la Tabla \ref{tab:exp2-dispersion}. Las dos coincidencias no contradicen el resto. Los
resultados posibles son múltiplos discretos de un ulp dentro de un rango acotado, así que con 50
muestras las repeticiones son esperables. La variabilidad tampoco es aleatoriedad del hardware: cada
permutación define un agrupamiento distinto y el resultado de cada uno es determinista, tanto que
fijar la semilla reproduce las 50 corridas.

La comparación con el orden decreciente pide cuidado en las dos direcciones. En mediana la
randomizada es mejor, $2.10 \times 10^{-14}$ contra $4.76 \times 10^{-14}$ para el mismo $N$, porque
mezclar rompe el patrón sistemático de absorción que produce recorrer los términos de mayor a menor.
Pero su máximo sobre la grilla grande es peor, $6.84 \times 10^{-14}$ contra
$5.03 \times 10^{-14}$. El orden decreciente es impreciso y predecible; el aleatorio es en mediana
algo más preciso y no reproducible sin fijar la semilla. Para un cálculo que debe dar el mismo valor dos
veces, la segunda propiedad pesa más que la primera.

\subsection{Representaciones equivalentes de una misma suma}
\label{subsec:dis-representaciones}

Las dos sucesiones de este experimento responden de manera opuesta a la misma transformación
algebraica, y esa oposición es lo que explica cuándo la cancelación importa.

Para $b_N$, las dos representaciones no producen el mismo flotante, cosa que era previsible porque
describen programas distintos según la sección \ref{subsec:telescopicas}, pero ninguna se degrada.
La representación telescópica resta $1/k - 1/(k+1)$, dos cantidades cuyo cociente tiende a 1, de modo que la
cancelación se agrava con $k$; y sin embargo la Figura \ref{fig:exp3-bn} muestra las dos curvas
contenidas por $\sqrt{N}\,u$ y sin tendencia creciente. La razón es que el término completo es de
orden $k^{-2}$ y se suma contra un acumulador que tiende a 1. Las cifras que la resta arruina son
cifras bajas de un sumando pequeño, y la acumulación las iba a descartar de todos modos. La
cancelación es un problema cuando el término cancelado domina el resultado, no cuando es una
corrección menor. La ventaja de la representación telescópica en la Tabla \ref{tab:exp3-resumen}, 97 contra
68 valores de $N$ con error nulo y un máximo algo menor, es de unos pocos ulp y no muestra tendencia
con $N$; no alcanza para preferir una representación sobre la otra.

Para $c_N$ el resultado es el contrario: $D_N$ no se anula en ninguno de los 1001 valores de $N$ y
crece casi seis órdenes de magnitud sobre la grilla, hasta $2.85 \times 10^{-11}$. Las dos
representaciones pasan de coincidir en 15 cifras a coincidir en 11. Acá el término cancelado es el
resultado mismo. El mecanismo es el descrito en la sección \ref{subsec:cancelacion} y merece
enunciarse al revés de como sugiere el nombre del fenómeno: por el lema de Sterbenz la resta
$\sqrt{k^{2}+1} - k$ es exacta, y el error lo aporta el redondeo previo de la raíz, de tamaño
absoluto del orden de $ku$, que la resta deja al frente de un resultado de tamaño $1/(2k)$. La
cancelación catastrófica no crea error, expone el que los operandos ya traían.

El acuerdo entre esa derivación y las mediciones es el resultado más ajustado del trabajo. La cota
$2k^{2}u$ de \eqref{eq:cota-termino} se obtuvo con aproximaciones de primer orden, y la Tabla
\ref{tab:exp3-terminos} muestra que los umbrales medidos caen entre 1.03 y 1.33 veces los
predichos, siempre del lado conservador. En el umbral final la coincidencia deja de ser aproximada:
la cota sitúa la pérdida completa de cifras significativas en $k = 1/\sqrt{2u} = 2^{26}$ y el primer
$k$ en que $\sqrt{k^{2}+1}$ devuelve exactamente $k$ es $67\,108\,864 = 2^{26}$. Un argumento de
orden de magnitud acierta un umbral de ocho cifras.

Dos advertencias sobre cómo leer esos datos. La primera es que $D_N$ mide la discrepancia entre dos
cálculos y no el error de ninguno de los dos, porque $c_N$ no tiene forma cerrada contra la cual
comparar (sección \ref{subsec:error-relativo}). Atribuir casi todo $D_N$ a la representación por
diferencia es razonable: la representación como cociente suma en el denominador y su error
relativo por término se mantiene del orden de $u$. Aun así, es una atribución apoyada en la teoría
y no una medición directa. La segunda es que $D_N$ queda muy por debajo de la suma de las
discrepancias de los términos individuales, que aparece en el panel (a) de la Figura
\ref{fig:exp3-cn}. Los errores de los términos tienen signos distintos y se compensan parcialmente, el mismo
comportamiento de paseo aleatorio que sostiene \eqref{eq:cota-raiz}. La suma de las cotas describe
el caso en que todos los errores se alinean, y presentarla como el error esperado sería tan
pesimista como usar $Nu$ para la suma acumulada.

Las dos formas cerradas de $b_N$ dieron el segundo resultado contraintuitivo. Por conteo de
operaciones, $1 - 1/(N+1)$ acumula dos redondeos y $N/(N+1)$ uno solo, así que \eqref{eq:modelo}
predice para la primera una cota del doble. Los errores máximos contra el valor racional exacto de
la Tabla \ref{tab:exp3-cerradas} confirman esa relación con precisión: $1.11 \times 10^{-16}$
contra $5.55 \times 10^{-17}$. Lo que no se sigue de ahí es que las dos expresiones den resultados
distintos en la práctica. Difieren en 4 de los 1001 valores de la grilla y en 19 de los primeros
$99\,999$, siempre por exactamente 1 ulp, y saturan en 1 con un solo $N$ de diferencia entre una y
otra. El redondeo adicional casi nunca llega a cambiar el flotante devuelto: para que lo cambie, el
valor exacto tiene que caer lo bastante cerca del punto medio entre dos representables, y eso ocurre
en una fracción pequeña de los casos.

De la comparación entre $c_N$ y las formas cerradas sale el punto general del experimento. Contar
operaciones acota el error pero no lo predice caso por caso: dos expresiones con cotas que difieren
en un factor 2 pueden devolver el mismo flotante casi siempre, mientras que dos expresiones con
cotas que difieren en un factor $k^{2}$ terminan sin ninguna cifra en común. La recomendación de
Goldberg (1991) de preferir la forma que evita restar cantidades próximas queda respaldada de manera
asimétrica: elegirla no empeoró ningún caso de este trabajo, y en el término de $c_N$ es, hacia
$k = 2^{26}$, la diferencia entre un error relativo del orden de $u$ y la pérdida completa de las
cifras significativas.

\subsection{Limitaciones}
\label{subsec:limitaciones}

El alcance de estas conclusiones está acotado por cuatro decisiones de diseño.

La referencia contra la cual se mide es, en la mayoría de los casos, un flotante. El error relativo
de \eqref{eq:error-relativo} se calcula contra $N/(N+1)$ evaluado en \texttt{binary64}, de modo que
el piso de lo medible es del orden de $u$ y los ceros de las Tablas \ref{tab:exp2-resumen} y
\ref{tab:exp3-resumen} significan coincidencia bit a bit con esa referencia, no exactitud. Solo la
comparación de las formas cerradas usa un patrón exacto, calculado con aritmética de fracciones. Un
trabajo posterior podría repetir las grillas grandes contra un valor de referencia en precisión
extendida y separar el error del algoritmo del error de la referencia.

Todas las mediciones corresponden a un único formato, \texttt{binary64}, sobre una máquina y un
intérprete. Los umbrales dependen del formato de manera predecible: en \texttt{binary32} el techo
$2^{53}$ del primer experimento pasaría a $2^{24}$ y el colapso de $c_N$ se adelantaría en más de
cuatro órdenes de magnitud en $k$. Eso es una extrapolación de las fórmulas y no una medición. Los
tiempos de la Tabla \ref{tab:exp1-costo} son los más sensibles al entorno: solo el cociente entre
las dos columnas es una magnitud estable entre corridas.

Las tres sucesiones estudiadas tienen términos positivos y decrecientes, y las conclusiones sobre el
orden de suma no se trasladan más allá de esa clase. En una serie con términos de signos alternados,
la cancelación entre términos domina el balance de error y acumular de menor a mayor módulo puede
dejar las cancelaciones grandes para el final, es decir, empeorar el resultado. La ventaja de
ordenar que se observa acá depende de que todos los términos empujen en la misma dirección.

Los rangos también acotan lo que se puede afirmar. El estudio de dispersión usa 50 repeticiones, que
alcanzan para el rango, la mediana y el desvío, pero no para describir las colas de la distribución.
Las grillas llegan a $N = 10^{6}$ para las sumas y a $k = 10^{8}$ para los términos de $c_N$; la
concordancia con $\sqrt{N}\,u$ está verificada dentro de esos límites y extrapolarla a tamaños
mucho mayores, donde el término $O(Nu^{2})$ de la cota de Kahan deja de ser despreciable, no está
respaldado por estos datos.
